\documentclass[10pt,a4paper]{amsart}
\usepackage[margin=19mm]{geometry}
\usepackage{amsmath,amssymb,mathtools}
\usepackage[scr=rsfs,cal=euler]{mathalfa}
\usepackage{xfrac}
\usepackage[unicode,hidelinks,bookmarks=false]{hyperref}
\newcommand{\ie}{i.e.\ }
\newcommand{\cf}{cf.\ }
\newcommand{\resp}{respectively\ }
\newcommand{\Subsection}{\subsection}
\providecommand{\nobreakdash}{\nobreak\hbox{-}}
\setlength{\emergencystretch}{2em}
\multlinegap0pt
\usepackage{bm}
\usepackage{bbm}
\usepackage{dsfont}
\usepackage{xspace}
\usepackage{cancel}
\usepackage{mathtools}
\usepackage{amsthm}
\makeatletter
\@ifundefined{lemma}{\newtheorem{lemma}{Lemma}}{}
\makeatother
\makeatletter
\newcommand{\E}{\mathbb{E}}
\newcommand{\Tr}{\operatorname{Tr}}
\newcommand{\abs}[1]{\left|#1\right|}
\def\bA{{\boldsymbol A}}
\def\bB{{\boldsymbol B}}
\def\bD{{\boldsymbol D}}
\def\bSigma{{\boldsymbol \Sigma}}
\def\bx{{\boldsymbol x}}
\def\msf{\mathsf}
\newcommand{\norm}[1]{\left\|#1\right\|}
\expandafter\ifx\csname proofoflemma\endcsname\relax
\newenvironment{proofoflemma}[1]{\par\noindent{\bf Proof of Lemma #1\ }}{\hfill\BlackBox\\[2mm]}
\fi
\expandafter\ifx\csname proofoftheorem\endcsname\relax
\newenvironment{proofoftheorem}[1]{\par\noindent{\bf Proof of Theorem #1\ }}{\hfill\BlackBox\\[2mm]}
\fi
\makeatother
\title[Universality of Kernel Random Matrices and Kernel Regression]{Universality of Kernel Random Matrices and Kernel Regression in the Quadratic Regime}
\author{Parthe Pandit, Zhichao Wang, Yizhe Zhu}
\date{Source: arXiv:2408.01062v2 (CC BY 4.0)}
\begin{document}
\maketitle

\noindent\emph{Source and licence.} Universality of Kernel Random Matrices and Kernel Regression in the Quadratic Regime. Version arXiv:2408.01062v2 (CC BY 4.0), distributed under the Creative Commons Attribution 4.0 International licence, as recorded in the arXiv licence metadata for this exact version and shown on the abstract page https://arxiv.org/abs/2408.01062v2. Full rights notice: licenses/arxiv-2408.01062-CC-BY-4.0.txt.\par\smallskip

\noindent\textit{Editorial scope.} Counted unit: the Lemma whose source label is \texttt{eq:centered\_x2} in the article (source0.tex lines 1584--1592, source label \texttt{eq:centered\_x2}), delivered together with the article's own complete original proof of it (source0.tex lines 1594--1681, one \texttt{proof} environment of 88 lines). No mathematics is written, repaired or re-derived: statement and proof are the article's own text line for line. Required context retained uncounted: eq:def\_reduced (lines 760--766). Every definition, display and label of those lines is retained unchanged. Omitted from this package and not redistributed: 1--759, 767--1583, 1593--1593, 1682--3354. Nothing inside a retained excerpt is omitted or altered: every retained body is delivered exactly as the article's lines are written, so no omission receipt of any kind accompanies this package. Citations: the retained excerpts carry no citation command at all, so no bibliography entry of the article is transcribed and no reference is redistributed. Reference adaptation: none was needed and none was made; every reference inside the delivered unit resolves inside it, so no unresolved reference or citation is produced. Printed slips kept literal: no printed word, symbol, hypothesis or number of the counted statement or of its proof was altered. Layout and class: the article's own class is \texttt{article} with packages including graphicx, booktabs, blindtext, bm, aux/jmlr2e, mathtools, enumerate, enumitem, tikz; the wrapper is the lane's pinned \texttt{amsart} editorial wrapper at 10pt on A4, which restates the theorem-like environments the retained text needs and adds the article's own macro definitions that the retained text needs (\texttt{\textbackslash E}, \texttt{\textbackslash Tr}, \texttt{\textbackslash abs}, \texttt{\textbackslash bA}, \texttt{\textbackslash bB}, \texttt{\textbackslash bD}, \texttt{\textbackslash bSigma}, \texttt{\textbackslash bx}, \texttt{\textbackslash msf}, \texttt{\textbackslash norm}), with extra wrapper packages (bm, bbm, dsfont, xspace, cancel, mathtools). The delivered numbering is the unit's own. Assets: no retained span contains a figure environment, a graphics-inclusion command, a PDF-page inclusion, a table or a TikZ picture; no image file is copied into this package, the package declares no asset dependency and no image file is redistributed with it. Licence: version arXiv:2408.01062v2 of this article is distributed under the Creative Commons Attribution 4.0 International licence, as recorded in the arXiv licence metadata for this exact version and shown on the abstract page https://arxiv.org/abs/2408.01062v2. A full rights notice is delivered at licenses/arxiv-2408.01062-CC-BY-4.0.txt. Characters: every retained excerpt is pure ASCII, so the delivered engine, XeLaTeX with its default Latin Modern text font, reports no missing character.
    \begin{equation}\label{eq:def_reduced}
    \bx_{i}^{(2)} (k,\ell)=
    \begin{cases}
        \sqrt{2}\bx_i(k)\bx_i(\ell) & k<\ell,\\
        \abs{\bx_i(k)}^2 & k=\ell.
    \end{cases}
    \end{equation}

\begin{lemma} Let $\bx\in \mathbb R^d$ be a random vector with independent entries and a diagonal covariance matrix $\bSigma$, where $\norm{\bSigma}\leq C$ for constant $C>0$. Assume each entry of $\bx$ has a zero mean and bounded 8th moments.   Let $\bx^{(2)}\in \mathbb R^{\binom{d+1}{2}}$ be a corresponding reduced tensor  vector  defined in \eqref{eq:def_reduced} and we define
      \begin{align}\label{eq:centered_x2}
      \overline\bx^{(2)}:=\bx^{(2)}-\E \bx^{(2)}.
      \end{align}
Then for any deterministic matrix $\bA$ with $\|\bA\|\leq 1$,
\begin{align}\label{eq:BZ_tensor}
    \E\left| {\overline\bx^{(2)}}^\top \bA\overline\bx^{(2)} -\Tr[\bA \bSigma^{(2)}]\right|^2=O(d^3).
     \end{align}
\end{lemma}

\begin{proof}
     We let $\bA=\bD+\bB \in \mathbb R^{\binom{d+1}{2}\times \binom{d+1}{2}}$, where $\bD$ is the diagonal part of $\bA$, and $\bB$ is the off-diagonal component of $\bA$.  Here the matrix $\bA$ is index by $\{(i,j): i\leq j, \quad i,j\in [d] \}$.
      To show \eqref{eq:BZ_tensor}, it suffices to bound the contribution from $\bD$ and $\bB$.



\textbf{ (i) Diagonal part.} Recall the definition of $\bx^{(2)}$ from \eqref{eq:def_reduced}. We have
\begin{align}
    &\E\left| {\overline\bx^{(2)}}^\top \bD\overline\bx^{(2)}-\Tr[\bD\bSigma^{(2)}]\right|^2\\
    =~&\E \left( \sum_{i<j}2(\bx_i^2 \bx_j^2-\bSigma_{ij,ij}^{(2)})\bA_{ij,ij} +\sum_{i} ((\bx_i^2-\bSigma_{ii})^2-\bSigma_{ii,ii}^{(2)}) \bA_{ii,ii}\right)^2 \notag \\
     \leq~& 4\sum_{i<j,k< l}|\bA_{ij,ij}\bA_{kl,kl}|\left|\E[ (\bx_{i}
    ^2\bx_j^2-\bSigma_{ij,ij}^{(2)})(\bx_{k}^2\bx_{l}^2-\bSigma_{kl,kl}^{(2)})]\right|\label{eq:D_off_term}\\
     &~~+\sum_{i,j} |\bA_{ii,ii}\bA_{jj,jj}| \left|\E[((\bx_i^2-\bSigma_{ii})^2-\bSigma_{ii,ii}^{(2)})((\bx_j^2-\bSigma_{jj})^2-\bSigma_{jj,jj}^{(2)})]\right|. \label{eq:D_diag_term}
\end{align}
    Since the $8$-th moments of $\bx_i$ are bounded for all $i\in [d]$, the contribution from \eqref{eq:D_diag_term} is at most $O(d)$. For \eqref{eq:D_off_term}, when $i,j,k,l$ are all distinct, by the diagonal assumption on $\bSigma$, $\bx_i,\bx_j,\bx_k,\bx_l$ are independent. We have
    $\E[ (\bx_{i}
    ^2\bx_j^2-\bSigma_{ij,ij}^{(2)})(\bx_{k}^2\bx_{l}^2-\bSigma_{kl,kl}^{(2)})]=0$.
Therefore, the nonzero contribution of \eqref{eq:D_off_term} only comes from indices $(i,j,k,l)$ that are not distinct. Since $\|\bD\|\leq \|\bA\|\leq 1$, we know the contribution with repeated indices  $(i,j,k,l)$ in \eqref{eq:D_off_term} is $O(d^3)$. Therefore, the total contribution from the diagonal part is $O(d^3)$.


\textbf{ (ii) Off-diagonal part.} We have the following expansion:
\begin{small}
\begin{align}
    &\E\left| {\overline\bx^{(2)}}^\top \bB\overline\bx^{(2)}-\Tr[\bB\bSigma^{(2)}]\right|^2 
    =\sum_{(i_1,i_2)\not=(i_3,i_4), (i_5,i_6)\not=(i_7,i_8)}\bA_{i_1i_2,i_3i_4}\bA_{i_5i_6,i_7i_8}\E[\overline\bx^{(2)}_{i_1i_2}\overline\bx^{(2)}_{i_3i_4}\overline\bx^{(2)}_{i_5i_6}\overline\bx^{(2)}_{i_7i_8}] \notag\\
    \lesssim & \sum_{(i_1,i_2)\not=(i_3,i_4), (i_5,i_6)\not=(i_7,i_8)} |\bA_{i_1i_2,i_3i_4}\bA_{i_5i_6,i_7i_8}|\label{eq:off_diag_8}\\
    &\cdot| \E[(\bx_{i_1}\bx_{i_2}-\bSigma_{i_1,i_2}\delta_{i_1,i_2})(\bx_{i_3}\bx_{i_4}-\bSigma_{i_3,i_4}\delta_{i_3,i_4})(\bx_{i_5}\bx_{i_6}-\bSigma_{i_5,i_6}\delta_{i_5,i_6})(\bx_{i_7}\bx_{i_8}-\bSigma_{i_7,i_8}\delta_{i_7,i_8})]|. 
\end{align}
\end{small}
For each index sequence $i_1,\dots, i_8$, to have a nonzero contribution in 
\begin{small}
\begin{align}\label{eq:i8}
\E[(\bx_{i_1}\bx_{i_2}-\bSigma_{i_1,i_2}\delta_{i_1,i_2})(\bx_{i_3}\bx_{i_4}-\bSigma_{i_3,i_4}\delta_{i_3,i_4})(\bx_{i_5}\bx_{i_6}-\bSigma_{i_5,i_6}\delta_{i_5,i_6})(\bx_{i_7}\bx_{i_8}-\bSigma_{i_7,i_8}\delta_{i_7,i_8})]
\end{align}
\end{small}
by the independence of the entries in $\bx$, there are at most 4 distinct values among $i_1,\dots, i_8$.  For sequences with at most 3 distinct indices, their total contribution in \eqref{eq:off_diag_8} is $O(d^3)$. Therefore, it suffices to estimate \eqref{eq:off_diag_8} when the contribution of index sequences with exactly 4 distinct indices satisfies $i_1\leq i_2,   i_3\leq i_4,  i_5\leq i_6,  i_7\leq i_8$. We have only the following cases depending on the number of distinct indices in $i_1,i_2,i_3,i_4$:
\begin{enumerate}
    \item Assume there are exactly 4 distinct indices in $i_1,\dots,i_4$. Then, to have a nonzero contribution, there is a perfect matching between $\{i_1,\dots, i_4\}$ and $\{i_5,\dots, i_8\}$. Using the inequality $2|\bA_{i_1i_2,i_3i_4}\bA_{i_5i_6,i_7i_8}|\leq |\bA_{i_1i_2,i_3i_4}|^2+|\bA_{i_5i_6,i_7i_8}|^2$, for an absolute constant $C$, the contribution is bounded by 
    \begin{align}
        C \left(\sum_{i_1< i_2,i_3< i_4} |\bA_{i_1i_2,i_3i_4}|^2\right)=C\|\bA\|_{\msf F}^2\leq C d^2 \|\bA\|^2 =O(d^2).
    \end{align}
    \item Assume there are exactly three distinct indices among $i_1,\dots, i_4$.  By symmetry, we only need to consider four subcases 
    \begin{itemize}
        \item (a) $i_1=i_2$, and $i_1,i_3,i_4$ are distinct. We can rewrite \eqref{eq:i8} as 
        \begin{align}\label{eq:i8a}
            \E[(\bx_{i_1}^2-\bSigma_{ii})\bx_{i_3}\bx_{i_4}(\bx_{i_5}\bx_{i_6}-\bSigma_{i_5,i_6}\delta_{i_5,i_6})(\bx_{i_7}\bx_{i_8}-\bSigma_{i_7,i_8}\delta_{i_7,i_8})].
        \end{align}
      Since there are exactly 4 distinct indices among $i_1,\dots,i_8$, and $i_1$ appears exactly twice,  $i_3,i_4,i_5,i_6,i_7,i_8$ must be distinct from $i_1$, which implies \eqref{eq:i8a} is equal to zero by independence.
    
      

        \item (b) $i_1=i_3$, and $i_1,i_2,i_4$ are distinct. We can rewrite \eqref{eq:i8} as 
        \begin{align}\label{eq:i8b}
            \E[\bx_{i_1}^2\bx_{i_2}\bx_{i_4}(\bx_{i_5}\bx_{i_6}-\bSigma_{i_5,i_6}\delta_{i_5,i_6})(\bx_{i_7}\bx_{i_8}-\bSigma_{i_7,i_8}\delta_{i_7,i_8})].
        \end{align}
   Note that if $i_5=i_6$ and $i_1,i_2,i_4,i_5$ are distinct, the expectation in \eqref{eq:i8b} is zero. By symmetry, we only need to consider $i_5=i_7, i_5=i_8$, or $i_5=i_2$.
     \begin{itemize}
         \item  (b.1) If $i_5=i_7$ and $i_1,i_2,i_4,i_5$ are distinct, we must have (i) $i_6=i_2$, $i_8=i_4$ or (ii) $i_6=i_4, i_8=i_2$. In case (i), we can bound \eqref{eq:off_diag_8} by 
        \begin{align}
            &\sum_{i_1\leq i_2, i_4, i_5} |\bA_{i_1i_2,i_1i_4}\bA_{i_5i_2, i_5,i_4}|\cdot \E[\bx_{i_1}^2\bx_{i_2}^2\bx_{i_4}^2\bx_{i_5}^2]\\
            \lesssim  &\sum_{i_1,i_2,i_4,i_5}\bA_{i_1i_2,i_1i_4}^2 +\sum_{i_1,i_2,i_4,i_5}\bA_{i_1i_2,i_1i_4}^2\lesssim d\|\bA\|_{\msf F}^2=O(d^3).
        \end{align}
        In case (ii), similarly, we can bound \eqref{eq:off_diag_8} by 
        \begin{align}
            &\sum_{i_1\leq i_2, i_4, i_5} |\bA_{i_1i_2,i_1i_4}\bA_{i_5i_4, i_5,i_2}|\cdot \E[\bx_{i_1}^2\bx_{i_2}^2\bx_{i_4}^2\bx_{i_5}^2]=O(d^3).
        \end{align}
        \item (b.2) If $i_5=i_8$, we must have (i) $i_6=i_2, i_7=i_4$ or (ii) $i_7=i_2, i_6=i_4$. In both cases, similar to case (b.1), the contribution is $O(d^3)$. 
        \item (b.3) If $i_5=i_2$, we must have (i) $i_7=i_4,i_8=i_6$ or (ii) $i_7=i_6, i_8=i_4$, and their contribution is $O(d^3)$.
     \end{itemize}
    
      
        
        \item (c) $i_2=i_4$, and $i_1,i_2,i_3$ are distinct. Like Case (b), its contribution is $O(d^3)$.
    \item (d) $i_1=i_4$ and $i_1,i_2,i_3$ are distinct. The same bound $O(d^3)$ holds.
    \end{itemize}
    


   
    \item  Assume there are exactly two distinct indices among $i_1,\dots, i_4$. We must have $i_1=i_2,i_3=i_4$, $i_1\not=i_3$ due to the constraint $(i_1,i_2)\not=(i_3,i_4)$.  In the same way, we must have $i_5=i_6, i_7=i_8, i_5\not=i_7$. 
Since there are 4 distinct indices among  $i_1,\dots,i_8$,   \eqref{eq:i8} becomes
  $ 
\E[(\bx_{i_1}^2-\bSigma_{i_1,i_1})(\bx_{i_3}^2-\bSigma_{i_3,i_3})(\bx_{i_5}^2-\bSigma_{i_5,i_5})(\bx_{i_7}^2-\bSigma_{i_7,i_7})]=0$.
    Therefore, the total contribution in this case is $0$.
\end{enumerate}
By the constraint $(i_1,i_2)\not=(i_3,i_4)$, there are at least 2 distinct indices among $i_1,\dots,i_4$. Therefore, we have discussed all three cases, and the total contribution for part (ii) is $O(d^3)$.
From the estimates in parts (i) and (ii) above, \eqref{eq:BZ_tensor} holds. 
\end{proof}

\end{document}
